\documentclass[11pt,letterpaper]{article}

\usepackage[utf8]{inputenc}
\usepackage[T1]{fontenc}
\usepackage[margin=1in]{geometry}
\usepackage{graphicx}
\usepackage{booktabs}
\usepackage{amsmath}
\usepackage{hyperref}
\usepackage{xcolor}
\usepackage{array}
\usepackage{float}
\usepackage{longtable}

\hypersetup{
  colorlinks=true,
  linkcolor=blue,
  urlcolor=blue,
  citecolor=blue,
  hypertexnames=false
}

\newcommand{\resultfigure}[4]{
  \begin{figure}[H]
  \centering
  \IfFileExists{#1}{
    \includegraphics[width=#2\textwidth]{#1}
  }{
    \fbox{\parbox{0.86\textwidth}{\centering
    Figura no encontrada: \texttt{\detokenize{#1}}.}}
  }
  \caption{#3}
  \label{#4}
  \end{figure}
}

\title{\textbf{Resilient Release Scheduling for the International Falcon Reservoir}\\
\large Challenge A -- Hackathon Latinoamerica 2026}
\author{Equipo de trabajo -- Falcon Reservoir}
\date{1 de julio de 2026}

\begin{document}

\maketitle

\begin{abstract}
Este reporte presenta una solucion reproducible para el Challenge A:
\textit{Resilient Release Scheduling for the International Falcon Reservoir}.
El objetivo es mejorar la resiliencia del almacenamiento del embalse durante
periodos de bajo volumen mediante una politica semanal de ajustes de liberacion.
La instancia oficial evaluada usa un horizonte de $T=26$ semanas y $L=5$
niveles discretos por semana, lo que produce una codificacion QUBO de 130
variables binarias y un espacio de busqueda de $5^{26}\approx 1.49\times10^{18}$
politicas posibles. La mejor politica nominal obtuvo
$\mathrm{SRS}=-0.145371$, con una mejora de
$\Delta\mathrm{SRS}=+0.010941$ respecto a la operacion historica. Ademas,
se realizo un analisis de robustez bajo siete escenarios hidrologicos; el
metodo QUBO+SA reoptimizado fue el mejor en el ranking agregado, con
\textit{robustez score} de $-0.456030$ y sin violaciones de factibilidad.
\end{abstract}

\section{Contexto del reto}

El International Falcon Reservoir forma parte del sistema binacional de la
cuenca del Rio Bravo / Rio Grande. El reto pregunta si una formulacion de
optimizacion basada en datos puede proponer calendarios de liberacion que
mejoren la resiliencia del almacenamiento durante periodos de escasez.

La variable de decision es el ajuste semanal de liberacion $u(t)$ aplicado
sobre la liberacion historica observada $R_{\mathrm{obs}}(t)$:

\begin{equation}
R_{\mathrm{opt}}(t) = R_{\mathrm{obs}}(t) + u(t).
\end{equation}

La dinamica simplificada del embalse queda:

\begin{equation}
S_{\mathrm{opt}}(t+1) =
S_{\mathrm{opt}}(t) + \Delta S_{\mathrm{obs}}(t) - u(t),
\end{equation}

donde $\Delta S_{\mathrm{obs}}(t)$ resume el cambio natural observado de
almacenamiento. Si $u(t)<0$, se libera menos agua que en la historia y se
conserva almacenamiento; si $u(t)>0$, se libera mas.

\section{Datos e instancia benchmark}

Los datos provienen del portal IBWC Water Data para las estaciones:

\begin{itemize}
  \item \textbf{08461200 -- International Falcon Reservoir}: almacenamiento,
  elevacion y cambio de almacenamiento.
  \item \textbf{08461300 -- Rio Grande Below Falcon Dam}: liberacion historica.
\end{itemize}

La instancia oficial media se configuro como:

\begin{table}[H]
\centering
\caption{Instancia oficial evaluada.}
\begin{tabular}{ll}
\toprule
Parametro & Valor \\
\midrule
Horizonte & $T=26$ semanas \\
Niveles discretos & $L=5$ \\
Variables binarias & $T\times L=130$ \\
Codificacion & One-hot por semana \\
Espacio de busqueda & $5^{26}\approx 1.49\times10^{18}$ politicas \\
Ventana de trabajo & 2024-01-01, 26 semanas \\
\bottomrule
\end{tabular}
\end{table}

La ventana inicia el 2024-01-01 porque en ese bloque de 26 semanas el sistema
atraviesa tres zonas operativas relevantes para el reto: almacenamiento bajo,
emergencia y condicion critica. Esto evita evaluar el algoritmo en un periodo
trivial y permite comparar las politicas cuando la presa se encuentra cerca de
los umbrales donde la resiliencia del almacenamiento es mas importante.

Los niveles discretos de ajuste son:

\begin{equation}
u(t) \in \{-2\Delta u,\,-\Delta u,\,0,\,\Delta u,\,2\Delta u\},
\qquad
\Delta u=0.25\,\mathrm{median}(R_{\mathrm{obs}}).
\end{equation}

En la corrida registrada, $\Delta u=3617.7255$ y el limite semanal es
$u_{\max}=2\Delta u$.

\section{Metrica oficial: Storage Resilience Score}

El Storage Resilience Score (SRS) se define como:

\begin{equation}
\mathrm{SRS} =
-\left(w_1 C_{\mathrm{crit}} + w_2 C_{\mathrm{dev}} +
w_3 C_{\mathrm{smooth}}\right).
\end{equation}

Por construccion, valores menos negativos son mejores. Los tres componentes
penalizan: almacenamiento bajo el umbral critico, desviaciones excesivas
respecto a la operacion historica y cambios bruscos entre semanas.

\begin{align}
C_{\mathrm{crit}} &=
\sum_{t=0}^{T}
\max(0, S_{\min} - S_{\mathrm{opt}}(t))^2, \\
C_{\mathrm{dev}} &=
\sum_{t=0}^{T-1} u(t)^2, \\
C_{\mathrm{smooth}} &=
\sum_{t=1}^{T-1} (u(t)-u(t-1))^2.
\end{align}

El umbral critico se calcula como:

\begin{equation}
S_{\min}=0.25S_{\max}.
\end{equation}

\section{Restricciones de factibilidad}

Todas las politicas reportadas se auditan contra las restricciones duras del
reto:

\begin{align}
R_{\mathrm{obs}}(t) + u(t) &\geq 0, \\
|u(t)| &\leq u_{\max}, \\
0 \leq S_{\mathrm{opt}}(t) &\leq S_{\max}, \\
\left|\sum_{t=0}^{T-1} u(t)\right| &\leq
0.10\sum_{t=0}^{T-1} R_{\mathrm{obs}}(t).
\end{align}

La cuarta restriccion evita que el modelo gane simplemente reteniendo agua en
todas las semanas; obliga a que la politica redistribuya temporalmente las
liberaciones dentro de un balance acumulado permitido.

\section{Metodos comparados}

Se compararon cinco rutas principales:

\begin{itemize}
  \item \textbf{Historico}: replay de la operacion observada, con $u(t)=0$.
  \item \textbf{Regla umbral}: reduce liberaciones cuando el almacenamiento cae
  por debajo de $S_{\min}$.
  \item \textbf{Genetico DEAP}: baseline clasico evolutivo sobre la politica
  discreta semanal.
  \item \textbf{QUBO + simulated annealing}: formulacion binaria one-hot con
  evaluacion final exacta de SRS y restricciones.
  \item \textbf{QCentroid}: ejecuciones separadas en modos \texttt{sa\_cpu},
  \texttt{sa\_gpu} y \texttt{qaoa} protegido.
\end{itemize}

La politica RNA se conserva como exploracion secundaria; no se reporta como
solucion final porque no mejora el baseline historico en esta corrida.

\section{Formulacion QUBO}

El ajuste semanal se codifica con variables binarias one-hot $x_{t,k}$:

\begin{equation}
u(t) = \Delta u \sum_k a_k x_{t,k},
\qquad
a=[-2,-1,0,1,2].
\end{equation}

El problema se lleva a una forma cuadratica:

\begin{equation}
\min_x x^\top Qx.
\end{equation}

La matriz $Q$ integra penalizaciones por desviacion, suavidad, almacenamiento
critico aproximado, one-hot, flujo no negativo y balance acumulado:

\begin{equation}
Q =
Q_{\mathrm{dev}}+
Q_{\mathrm{smooth}}+
Q_{\mathrm{crit}}+
Q_{\mathrm{onehot}}+
Q_{\mathrm{flow}}+
Q_{\mathrm{balance}}.
\end{equation}

El termino $Q_{\mathrm{crit}}$ usa una aproximacion cuadratica para permitir
una forma QUBO. Por eso, la seleccion final no se toma solo por energia QUBO:
cada candidato se decodifica y se reevalua con el SRS oficial y la auditoria
de factibilidad.

\section{Resultados nominales}

La Tabla \ref{tab:nominal} resume la corrida principal. El mejor SRS nominal
lo alcanzan tanto el genetico DEAP como QUBO+SA. La mejora respecto al replay
historico es de $+0.010941$.

\begin{table}[H]
\centering
\caption{Resultados nominales para $T=26$, $L=5$.}
\label{tab:nominal}
\begin{tabular}{lrrr}
\toprule
Metodo & SRS & $\Delta$SRS vs historico & Tiempo (s) \\
\midrule
Historico ($u=0$) & -0.156312 & 0.000000 & -- \\
Regla umbral & -0.149991 & +0.006321 & -- \\
Genetico DEAP & -0.145371 & +0.010941 & 7.89 \\
QUBO + SA & -0.145371 & +0.010941 & 68.60 \\
RNA exploratoria & -0.156776 & -0.000465 & -- \\
\bottomrule
\end{tabular}
\end{table}

La politica QUBO+SA seleccionada reduce la liberacion durante las primeras
11 semanas con $u(t)=-\Delta u$ y mantiene $u(t)=0$ en las 15 semanas
restantes. Sus indicadores principales son:

\begin{table}[H]
\centering
\caption{Resumen de la politica QUBO+SA seleccionada.}
\begin{tabular}{lr}
\toprule
Indicador & Valor \\
\midrule
SRS & -0.145371 \\
$\Delta$SRS vs historico & +0.010941 \\
Semanas con reduccion de liberacion & 11 \\
Semanas sin ajuste & 15 \\
$\sum_t u(t)$ & -39794.9805 \\
Almacenamiento minimo simulado & 348251.0805 \\
Almacenamiento final simulado & 807176.0805 \\
Violaciones de restricciones & 0 \\
\bottomrule
\end{tabular}
\end{table}

\section{Resultados QCentroid}

La Tabla \ref{tab:qcentroid} compara las corridas exportadas desde los
notebooks de QCentroid. Los tres modos produjeron la misma politica final y
el mismo SRS. El modo \texttt{qaoa} debe interpretarse con cuidado: para la
instancia de 130 qubits, la ejecucion protegida registro \texttt{fallback}
desde QAOA hacia \texttt{sa\_gpu}; por lo tanto, se reporta como ruta
QAOA disponible, no como demostracion de ventaja cuantica pura.

\begin{table}[H]
\centering
\caption{Corridas QCentroid exportadas.}
\label{tab:qcentroid}
\begin{tabular}{lrrl}
\toprule
Modo & SRS & Tiempo (s) & Observacion \\
\midrule
\texttt{sa\_cpu} & -0.145371 & 97.46 & CPU reproducible \\
\texttt{sa\_gpu} & -0.145371 & 4350.77 & $128$ restarts, $500000$ iteraciones \\
\texttt{qaoa} protegido & -0.145371 & 4370.69 & fallback desde QAOA; 130 qubits solicitados \\
\bottomrule
\end{tabular}
\end{table}

Esta lectura es importante para la defensa tecnica: la contribucion cuantica
esta en la formulacion QUBO/Ising y en la compatibilidad con una ruta QAOA,
pero el resultado oficial estable para la instancia media se obtiene con
annealing hibrido/clasico sobre la formulacion binaria.

\section{Validacion QAOA reducida y benchmark de escalabilidad}

Adicionalmente se reviso el notebook
\texttt{Versiones/reservoir\_falcon\_benchmark.ipynb}. Su principal aporte es
separar dos preguntas que conviene no mezclar durante la presentacion:

\begin{itemize}
  \item \textbf{Correctitud de la ruta QAOA}: se valido una ventana real reducida
  de $T=6$, $L=3$, equivalente a 18 qubits. En ese caso, QAOA statevector alcanzo
  el mismo minimo QUBO que una busqueda exhaustiva.
  \item \textbf{Escalabilidad del caso oficial}: la instancia oficial
  $T=26$, $L=5$ requiere 130 qubits y representa
  $1.49\times10^{18}$ politicas posibles. Esto justifica reportar el caso
  oficial con QUBO+SA y dejar QAOA completo como ruta experimental.
\end{itemize}

\begin{table}[H]
\centering
\caption{Aportes del benchmark QAOA/escalabilidad agregado.}
\begin{tabular}{lrrrl}
\toprule
Instancia & $T$ & $L$ & Qubits & Resultado \\
\midrule
Validacion QAOA & 6 & 3 & 18 & QAOA alcanzo el minimo QUBO exacto \\
Small & 12 & 3 & 36 & SA validado contra brute force \\
Medium oficial & 26 & 5 & 130 & QUBO+SA estable para entrega \\
\bottomrule
\end{tabular}
\end{table}

\section{Analisis de robustez}

Para evaluar si la politica solo funciona en el caso nominal o conserva valor
bajo incertidumbre, se construyeron siete escenarios:

\begin{itemize}
  \item nominal;
  \item sequia 10\%;
  \item sequia 20\%;
  \item ruido semanal suave;
  \item ruido semanal fuerte;
  \item bloque critico de cinco semanas;
  \item desfase temporal del cambio de almacenamiento.
\end{itemize}

Se compararon cuatro politicas: historico, regla umbral, QUBO+SA nominal y
QUBO+SA reoptimizado por escenario. El ranking agregado se muestra en la
Tabla \ref{tab:robustez}.

\begin{table}[H]
\centering
\caption{Ranking de robustez bajo escenarios perturbados.}
\label{tab:robustez}
\begin{tabular}{lrrrr}
\toprule
Metodo & SRS prom. & SRS peor & Factibilidad & Score robustez \\
\midrule
QUBO+SA reoptimizado & -0.162869 & -0.217459 & 1.000 & -0.456030 \\
QUBO+SA nominal & -0.163411 & -0.219375 & 1.000 & -0.457947 \\
Regla umbral & -0.167510 & -0.219708 & 0.857 & -0.468279 \\
Historico & -0.176004 & -0.239564 & 1.000 & -0.488136 \\
\bottomrule
\end{tabular}
\end{table}

El resultado refuerza dos puntos. Primero, QUBO+SA no solo mejora el caso
nominal: tambien domina el promedio y el peor caso entre las politicas
evaluadas. Segundo, reoptimizar ante cada escenario produce una mejora
pequena pero consistente sobre usar la politica nominal sin ajuste.

\section{Figuras generadas}

\resultfigure{resultados/simulacion_hibrida.png}{0.96}
{Trayectorias de almacenamiento y ajustes semanales para la corrida principal.}
{fig:simulacion}

\resultfigure{resultados/resultados_sa_gpu/qubo_matrix.png}{0.76}
{Matriz QUBO generada para la instancia $T=26$, $L=5$.}
{fig:qubo-matrix}

\resultfigure{resultados/resultados_sa_gpu/qubo_components.png}{0.86}
{Componentes de la formulacion QUBO usada por la ruta QCentroid SA GPU.}
{fig:qubo-components}

\resultfigure{resultados/resultados_robustez/ranking_robustez.png}{0.82}
{Ranking agregado de robustez bajo escenarios hidrologicos perturbados.}
{fig:ranking}

\resultfigure{resultados/resultados_robustez/heatmap_semanas_criticas.png}{0.92}
{Semanas criticas por metodo y escenario.}
{fig:heatmap}

\resultfigure{resultados/resultados_robustez/boxplot_srs_robustez.png}{0.86}
{Distribucion del SRS por metodo en los escenarios de robustez.}
{fig:boxplot}

\resultfigure{resultados/resultados_benchmark_qaoa/qaoa_validation_18q.png}{0.68}
{Validacion QAOA en una ventana real reducida de 18 qubits; QAOA alcanza el minimo QUBO exacto.}
{fig:qaoa-validation}

\resultfigure{resultados/resultados_benchmark_qaoa/scaling_search_runtime.png}{0.92}
{Escalabilidad: crecimiento del espacio de busqueda y tiempo de solver quantum-inspired.}
{fig:scaling-qaoa}

\resultfigure{resultados/resultados_benchmark_qaoa/cost_breakdown_srs.png}{0.82}
{Desglose de contribuciones ponderadas al costo $-\mathrm{SRS}$ para el caso oficial.}
{fig:cost-breakdown}

\section{Escalabilidad}

El crecimiento combinatorio justifica el uso de una formulacion binaria y de
metodos hibridos. La busqueda exhaustiva deja de ser viable desde la instancia
media:

\begin{table}[H]
\centering
\caption{Crecimiento del espacio de busqueda.}
\begin{tabular}{lrrr}
\toprule
Instancia & Semanas $T$ & Niveles $L$ & Politicas posibles \\
\midrule
Small & 12 & 3 & $3^{12}=531441$ \\
Medium oficial & 26 & 5 & $5^{26}\approx 1.49\times 10^{18}$ \\
Large & 52 & 5 & $5^{52}\approx 2.22\times 10^{36}$ \\
\bottomrule
\end{tabular}
\end{table}

La formulacion QUBO escala linealmente en numero de variables binarias
respecto a $T\times L$, aunque la densidad de acoplamientos y la dificultad de
optimizacion crecen por las penalizaciones de suavidad, balance y trayectoria
de almacenamiento.

\section{Limitaciones}

El modelo de embalse utilizado es deliberadamente simplificado, como lo indica
el reto. No sustituye reglas operativas oficiales ni incorpora todos los
factores legales, agricolas, urbanos o politicos del sistema binacional.

Tambien se debe evitar sobreinterpretar la ruta QAOA: la instancia media
requiere 130 qubits en la codificacion one-hot, por lo que la corrida
registrada activo una proteccion y uso \texttt{sa\_gpu} como solucion estable.
El valor tecnico de esa ruta es mostrar que el problema esta correctamente
codificado en forma QUBO/Ising y que puede conectarse a backends cuanticos o
hibridos, no afirmar una ventaja cuantica experimental.

\section{Uso de herramientas de apoyo}

Durante el desarrollo se usaron herramientas de apoyo para revision de codigo,
organizacion de documentacion, depuracion y redaccion tecnica. El equipo definio
la formulacion del problema, selecciono los datos hidrologicos, ejecuto los
notebooks y scripts de optimizacion, verifico los resultados numericos y es
responsable de la interpretacion final. Las decisiones tecnicas centrales
--metrica SRS, restricciones, discretizacion, QUBO, comparacion contra baselines
y analisis de robustez-- se reportan con archivos reproducibles dentro del
repositorio.

\section{Conclusiones}

La solucion construida cumple los elementos centrales solicitados por el
Challenge A:

\begin{itemize}
  \item discretiza ajustes semanales de liberacion;
  \item codifica la politica como variables binarias one-hot;
  \item deriva y ejecuta una formulacion QUBO de 130 variables;
  \item compara contra historico, regla umbral y genetico DEAP;
  \item reporta $\Delta$SRS y tiempos de computo;
  \item audita restricciones de factibilidad;
  \item documenta modos QCentroid \texttt{sa\_cpu}, \texttt{sa\_gpu} y
  \texttt{qaoa};
  \item agrega una evaluacion de robustez bajo escenarios perturbados.
\end{itemize}

El resultado principal es una mejora nominal de
$\Delta\mathrm{SRS}=+0.010941$ frente al replay historico y una politica QUBO+SA
factible. En robustez, QUBO+SA reoptimizado obtuvo el mejor ranking agregado
entre los metodos evaluados.

\section*{Archivos de reproduccion}

\begin{itemize}
  \item \texttt{resultados/benchmark\_delta\_srs.csv}
  \item \texttt{resultados/resultados\_sa\_cpu/benchmark\_delta\_srs (1).csv}
  \item \texttt{resultados/resultados\_sa\_gpu/benchmark\_delta\_srs.csv}
  \item \texttt{resultados/resultados\_QAOA/benchmark\_delta\_srs.csv}
  \item \texttt{resultados/resultados\_robustez/resumen\_robustez.csv}
  \item \texttt{resultados/resultados\_robustez/detalle\_escenarios.csv}
  \item \texttt{FalconChallenge\_QCentroid\_sa\_cpu.ipynb}
  \item \texttt{FalconChallenge\_QCentroid\_sa\_gpu.ipynb}
  \item \texttt{FalconChallenge\_QCentroid.ipynb}
\end{itemize}

\section*{Referencias}

\begingroup
\sloppy
\begin{itemize}
  \item United Nations, Sustainable Development Goal 6: \url{https://sdgs.un.org/goals/goal6}.
  \item United Nations, Sustainable Development Goal 13: \url{https://sdgs.un.org/goals/goal13}.
  \item E. Farhi, J. Goldstone, S. Gutmann, \textit{A Quantum Approximate Optimization Algorithm}, arXiv:1411.4028, 2014. \url{https://arxiv.org/abs/1411.4028}.
  \item A. Lucas, \textit{Ising formulations of many NP problems}, Frontiers in Physics, 2014. \url{https://www.frontiersin.org/journals/physics/articles/10.3389/fphy.2014.00005/full}.
  \item International Boundary and Water Commission Water Data Portal: \url{https://waterdata.ibwc.gov}.
\end{itemize}
\endgroup

\end{document}
